\documentclass[man,12pt,floatsintext,donotrepeattitle]{apa7}

\usepackage[american]{babel}
\usepackage{csquotes}
\usepackage{amsmath,amssymb}
\usepackage{booktabs}
\usepackage{xcolor}
\usepackage[most]{tcolorbox}
\usepackage[style=apa,backend=biber]{biblatex}
\usepackage{xr}
\externaldocument{draft6}
\addbibresource{refs.bib}
\hypersetup{
  hidelinks,
  pdftitle={Supplement 2: Definitions, Theorems and Worked Examples},
  pdfauthor={Joseph Corneli},
  pdfsubject={Formal apparatus accompanying Rule-Rewriting Cellular Automata
    and the Edge of Chaos}
}

\newcommand{\sig}{\sigma}
\newcommand{\bit}[1]{\mathrm{bit}[#1]}
% Cross-reference a finding box by its label (tcolorbox auto counter supplies the number).
\newcommand{\findref}[1]{Box~\ref{#1}}


\title{Supplement 2: Definitions, Theorems and Worked Examples}
\shorttitle{Definitions, Theorems and Worked Examples}
\authorsnames{Joseph Corneli}
\authorsaffiliations{{Oxford Brookes University, Oxford OX3 0BP, United Kingdom},
  {Hyperreal Enterprises Ltd, United Kingdom}}
\note{Supplement to \emph{Rule-Rewriting Cellular Automata and the Edge of
Chaos}. Correspondence: \href{mailto:jcorneli@brookes.ac.uk}
{jcorneli@brookes.ac.uk}.}

\abstract{%
This notebook collects the twelve empirical findings referenced from the main
paper. The numbering is stable: ``Supplement 1, Box \(n\)'' in the paper points
to Box \(n\) here. Mathematical definitions, protocols, figures, and
interpretation remain in the main paper; this supplement isolates the compact
claims that those materials support. A parallel Clay rendering provides a
public-facing literate-notebook view from the same canonical source.}

% --- typed-box environments, carried from draft6.tex (never folded here) ---
\newif\ifpaperboxfold
\tcbset{%
  fold/.code={%
    \paperboxfoldfalse
    \pgfkeysalso{top=0pt,bottom=0pt,boxsep=0pt}},
  unfold/.code={\paperboxfoldfalse}}
\newcommand{\foldmark}{%
  \ifpaperboxfold\space{\normalfont\scriptsize[folded]}\fi}

\newtcolorbox[auto counter]{workedexampleframe}[2][]{%
  enhanced,
  breakable,
  colback=green!6,
  colframe=green!45!black,
  boxrule=0.7pt,
  arc=1.5mm,
  left=2.5mm,
  right=2.5mm,
  top=1.5mm,
  bottom=1.5mm,
  fonttitle=\bfseries,
  title={Worked example~\thetcbcounter: #2\foldmark},
  #1}
\NewDocumentEnvironment{workedexample}{O{} m +b}{%
  \paperboxfoldfalse
  \begin{workedexampleframe}[#1]{#2}%
    \ifpaperboxfold\else#3\fi
  \end{workedexampleframe}}{}

% --- typed boxes for compression -------------------------------------------
% Each typed box accepts `fold` in its option list. Folding preserves its
% counter, title, label, and cross-references while suppressing only the body.
\newtcolorbox[auto counter]{theoremframe}[1][]{%
  enhanced, breakable, colback=blue!5, colframe=blue!55!black,
  boxrule=0.7pt, arc=1.5mm, left=2.5mm, right=2.5mm,
  top=1.5mm, bottom=1.5mm, fonttitle=\bfseries,
  title={Theorem~\thetcbcounter\foldmark}, #1}
\NewDocumentEnvironment{theorem}{O{} +b}{%
  \paperboxfoldfalse
  \begin{theoremframe}[#1]%
    \ifpaperboxfold\else#2\fi
  \end{theoremframe}}{}

\newtcolorbox[auto counter]{definitionframe}[2][]{%
  enhanced, breakable, colback=violet!6, colframe=violet!55!black,
  boxrule=0.7pt, arc=1.5mm, left=2.5mm, right=2.5mm,
  top=1.5mm, bottom=1.5mm, fonttitle=\bfseries,
  title={Definition~\thetcbcounter: #2\foldmark}, #1}
\NewDocumentEnvironment{definition}{O{} m +b}{%
  \paperboxfoldfalse
  \begin{definitionframe}[#1]{#2}%
    \ifpaperboxfold\else#3\fi
  \end{definitionframe}}{}

% Corollaries auto-boxed like theorems (blue). Empirical findings are amber:
% observed (not proved) claims live here, the validated home for rewritten results.
\newtcolorbox[auto counter]{corollaryframe}[1][]{%
  enhanced, breakable, colback=blue!5, colframe=blue!55!black,
  boxrule=0.7pt, arc=1.5mm, left=2.5mm, right=2.5mm,
  top=1.5mm, bottom=1.5mm, fonttitle=\bfseries,
  title={Corollary~\thetcbcounter\foldmark}, #1}
\NewDocumentEnvironment{corollary}{O{} +b}{%
  \paperboxfoldfalse
  \begin{corollaryframe}[#1]%
    \ifpaperboxfold\else#2\fi
  \end{corollaryframe}}{}

\newtcolorbox[auto counter]{propositionframe}[2][]{%
  enhanced, breakable, colback=blue!5, colframe=blue!55!black,
  boxrule=0.7pt, arc=1.5mm, left=2.5mm, right=2.5mm,
  top=1.5mm, bottom=1.5mm, fonttitle=\bfseries,
  title={Proposition~\thetcbcounter: #2\foldmark}, #1}
\NewDocumentEnvironment{proposition}{O{} m +b}{%
  \paperboxfoldfalse
  \begin{propositionframe}[#1]{#2}%
    \ifpaperboxfold\else#3\fi
  \end{propositionframe}}{}

\newtcolorbox[auto counter]{empiricalfindingframe}[2][]{%
  enhanced, breakable, colback=orange!7, colframe=orange!65!black,
  boxrule=0.7pt, arc=1.5mm, left=2.5mm, right=2.5mm, top=1.5mm, bottom=1.5mm,
  fonttitle=\bfseries,
  title={Empirical finding~\thetcbcounter: #2\foldmark}, #1}
\NewDocumentEnvironment{empiricalfinding}{O{} m +b}{%
  \paperboxfoldfalse
  \begin{empiricalfindingframe}[#1]{#2}%
    \ifpaperboxfold\else#3\fi
  \end{empiricalfindingframe}}{}

% Empirical findings live canonically in the literate notebook submitted as
% Supplement 1. Keep only stable, compact pointers in the narrative.
\newcommand{\suppfindingref}[1]{Supplement~1, Box~#1}
\newcommand{\suppfinding}[2]{%
  \par\smallskip\noindent\emph{\suppfindingref{#1}: #2.}\par\smallskip}

% Observations (teal): consequences/remarks pulled out of definitions and prose.
\newtcolorbox[auto counter]{observationframe}[2][]{%
  enhanced, breakable, colback=teal!7, colframe=teal!55!black,
  boxrule=0.7pt, arc=1.5mm, left=2.5mm, right=2.5mm, top=1.5mm, bottom=1.5mm,
  fonttitle=\bfseries, title={Observation~\thetcbcounter: #2\foldmark}, #1}
\NewDocumentEnvironment{observation}{O{} m +b}{%
  \paperboxfoldfalse
  \begin{observationframe}[#1]{#2}%
    \ifpaperboxfold\else#3\fi
  \end{observationframe}}{}


\newcommand{\secref}[1]{\emph{\nameref{#1}}}

\begin{document}

\part*{Supplement 2: Definitions, Theorems and Worked Examples}

This supplement holds the formal apparatus of the main paper: the definitions
that fix the object, the two theorems and their corollaries that classify it,
and the worked examples that carry a reader through one member of the family by
hand. The main paper states each result in prose where it is used and points
here for the statement, the proof, and the arithmetic. Numbering matches the
order of first use in the paper.


\begin{definition}[label={def:writing}]{Writings, elementary writes, and propagators}
A \emph{writing} is any function $w\colon\{0,\dots,7\}\to\{0,\dots,7\}$ on the
truth-table positions. Its elementary update at source position $p$ is
\[
(U_{w,p}g)(q)=
\begin{cases}
\lnot g(p),&q=w(p),\\
g(q),&q\ne w(p).
\end{cases}
\]
Thus one update complements source bit $p$ and writes it to target $w(p)$. If
$w$ is a bijection---a permutation $\sig\in S_8$---all eight writes can be
performed simultaneously without collisions, giving the
\emph{permutation-propagator} $P_\sig$ in Equation~\eqref{eq:prop}. Its fixed
bytes are exactly the bytes fixed by every $U_{\sig,p}$. For a non-injective
writing, individual updates $U_{w,p}$ remain well defined, but a simultaneous
$P_w$ requires a collision-resolution convention; such a writing also has
\emph{gaps}, or positions targeted by no source.
\end{definition}

\begin{observation}[label={obs:core}]{Permutations form a small bijective core}
Permutations are only the bijective core of a far larger class: there are
$8!=40{,}320$ permutations but $8^{8}=16{,}777{,}216$ writings, so the
permutations are $0.24\%$ of the whole. The cycle-structure algebra of
\secref{sec:classification} is available only on this core. The historical
operator $[0\,0\,1\,2\,3\,4\,5\,6]$ (Figure~\ref{fig:fig2pair}) is itself a
non-bijective writing---bits $0$ and $1$ both write to position $0$, and
position $7$ is never written---and its behaviour ranges widely between its two
standard-order forms.
\end{observation}

\begin{workedexample}[label={box:propagator}]{A simultaneous propagator by hand}
\small
Write an ECA rule in Wolfram's standard neighbourhood order:
\begin{center}
\setlength{\tabcolsep}{3.2pt}
\begin{tabular}{c@{\quad}*{8}{c}}
position & $0$ & $1$ & $2$ & $3$ & $4$ & $5$ & $6$ & $7$ \\
neighbourhood & $111$ & $110$ & $101$ & $100$ & $011$ & $010$ & $001$ & $000$ \\
\midrule
Rule 110 & $0$ & $1$ & $1$ & $0$ & $1$ & $1$ & $1$ & $0$ \\
\end{tabular}
\end{center}
Thus Rule~110 is the byte $01101110$. Let $\sig(p)=p+2\pmod 8$ on the
displayed position indices: each source bit is complemented and written two
positions further to the right, wrapping at the end.
Carrying out the eight writes gives
\[
P_{+2}(01101110)=01100100,
\]
which is Rule~100. This changes what the cell does to its black-and-white
phenotype. For example, neighbourhood $011$ maps to $1$ under Rule~110 but
to $0$ under Rule~100. The propagator has not directly flipped a phenotype
cell; it has changed the truth table that the cell will use.
\end{workedexample}

\begin{definition}[label={def:genopheno}]{Genotype and phenotype}
Each cell carries an eight-bit rule. The field of rules---the
\emph{genotype}---governs a field of black-and-white states---the
\emph{phenotype}---and neighbouring rules also interact. One complete local
update (Box~\ref{box:metaca-step}) evolves both layers; this is the computation
behind the space--time diagrams in Figure~\ref{fig:regimes}.
\end{definition}

\begin{workedexample}[label={box:metaca-step}]{One site, one complete MetaCA step}
\small
At time $t$, site $x$ holds an eight-bit rule $g_x(t)$ and a phenotype bit
$X_x(t)$. One ordinary coupled step has three parts.
\begin{enumerate}
\item \textbf{Express the current rule.} The phenotype neighbourhood is
evaluated by the rule at its centre. With numerical neighbourhood value
$n_x(t)=4X_{x-1}(t)+2X_x(t)+X_{x+1}(t)$, its displayed position is $7-n_x(t)$:
\[
X_x(t+1)=g_x(t)\bigl(7-n_x(t)\bigr).
\]
\item \textbf{Blend neighbouring rules.} For each truth-table position $j$,
inspect the three genotype bits
$\bigl(g_{x-1}(j),g_x(j),g_{x+1}(j)\bigr)$. If the two outer bits agree, copy
their value. Thus $(0,1,0)$ blends to $0$. If they disagree, the centre rule
decides by treating the triple as one of its ECA neighbourhoods, converting
its numerical value $n$ to displayed position $7-n$ as above. The eight
coordinate-wise decisions form an intermediate rule
\[
h_x=B\bigl(g_{x-1}(t),g_x(t),g_{x+1}(t)\bigr).
\]
\item \textbf{Write one propagated bit.} Draw a source position
$K_x(t)\in\{0,\ldots,7\}$ uniformly and apply its elementary write to the
intermediate byte:
\[
g_x(t+1)=U_{\sig,K_x(t)}(h_x).
\]
\end{enumerate}
The ordinary spatial runs make one such elementary write per cell and
generation. The historical Figure~8 variant made two. By contrast,
$P_\sig$ denotes the simultaneous eight-write map classified in
\secref{sec:classification}; its fixed bytes are precisely the bytes left
unchanged by every possible elementary write.
\end{workedexample}

\begin{theorem}[label={thm:parity}]
$P_\sig$ has a fixed point (a fixed rule) if and only if every cycle of
$\sig$ has even length. Since a one-cycle has odd length, this already
requires $\sig$ to be fixed-point-free.
\end{theorem}

\begin{workedexample}[label={box:even-cycle}]{Even and odd cycles}
\small
For offset $+2$ on eight truth-table positions,
\[
\sig=(0\;2\;4\;6)(1\;3\;5\;7).
\]
A fixed rule must change colour at every arrow because
$g(\sig k)=\lnot g(k)$. Around either four-cycle it can therefore read
$0101$ or $1010$: after four negations it returns consistently to its
starting bit. The two four-cycles are coloured independently, so two binary
choices fix all eight truth-table entries $\bit{0},\dots,\bit{7}$. Read back as
an ECA rule number---those eight bits in Wolfram neighbourhood order (the table
of \secref{sec:object}) taken as an integer in $0$--$255$---the four colourings
are the fixed rules
\[
2^2=4:\quad
51,\ 102,\ 153,\ 204,
\]
each with four zero-bits and four one-bits. An odd cycle admits no such
colouring: after an odd number of negations the bit returns to its starting
position demanding its own opposite, so no fixed rule exists. Admitting a fixed
rule and settling onto one are distinct. An operator carrying an odd cycle has no
fixed rule for its field to collapse onto, and its field correspondingly stays
diverse (the any-odd reduced rotate$+2$, Figure~\ref{fig:parity}, right); whether
a fixed-point-bearing operator collapses onto one of its fixed rules is a
separate, dynamical question (\secref{sec:interrupter}).
\end{workedexample}

\begin{corollary}[label={cor:count}]
An all-even $\sig$ with $k$ cycles has exactly $2^k$ fixed bytes; any $\sig$ with
an odd cycle has none.
\end{corollary}

\begin{definition}[label={def:lambda}]{Langton's activity $\lambda$}
$\lambda$ measures a rule's \emph{activity}: the fraction of its neighbourhood
entries that map to the non-quiescent
state~\parencite{langton1990computation}---for an eight-bit ECA table, the
Hamming weight over eight, $\lambda = \lVert r \rVert_1 / 8$, from $0$ (every
neighbourhood dies) through $1/2$ to $1$. A rule with $\lambda = 1/2$ is
\emph{balanced}.
\end{definition}

\begin{observation}[label={obs:lambda}]{Balance is a truth-table property}
Langton's thesis was that complex behaviour concentrates near a critical
$\lambda$, shown most clearly for larger tables ($K=4$, $N=5$); for the
two-symbol, three-neighbour ECAs Langton is explicit that $\lambda$ only roughly
correlates with dynamics, and later work is
cautious~\parencite{mitchell1993revisiting}. Theorem~\ref{thm:lambda} below
concerns the \emph{value} $\lambda = 1/2$ itself---balance---which we keep
separate from any empirical claim about what it predicts dynamically.
\end{observation}

\begin{theorem}[label={thm:lambda}]
Every fixed point of every operator in the family has exactly four one-bits;
equivalently, Langton's activity parameter $\lambda = 4/8 = 1/2$.
\end{theorem}

\begin{definition}[label={def:aliveness}]{Census protocol and aliveness}
For each of the $40{,}320$ permutations, we ran seeds $0$--$2$ on a width-$60$
line for $100$ generations. Initial rule bytes are independent uniform draws
from $0$--$255$, initial phenotype bits are independent fair binary draws, and
both layers have fixed zero boundaries. Each genotype update uses the
neighbour blend and one uniformly sampled elementary write described in
Box~\ref{box:metaca-step}. A layer's \emph{change rate} is the mean fraction of
cells differing between consecutive rows over the final $40$ generations,
averaged over the three seeds. The layer is \emph{alive} when this rate is at
least $0.10$ and \emph{dead} otherwise. Scoring genotype and phenotype
separately gives four regimes: live/live, live/dead, dead/live, and dead/dead.
\end{definition}

\begin{workedexample}[label={box:survey},breakable=false]{Reading the rotation survey}
\small
The fields of Supplement~3, Figure~1 evolve under the \emph{spatially coupled}
dynamics of \secref{sec:interrupter}---stochastic elementary writes together
with the coupling between cells---not $P_\sig$ applied cell-by-cell. This is
what the survey turns on. Without the blend, Figure~\ref{fig:parity} already
shows rotate$+1$ remaining active over the observed horizon while rotate$+2$
settles. With neighbour blending, the finite-horizon split persists and widens:
the single 8-cycles---$\gcd(\text{offset},8)=1$, offset $\pm 1$ and $\pm 3$---
sustain a high-diversity field, while offset $\pm 2$ ($\gcd 2$) and $\pm 4$
($\gcd 4$) collapse; that split is what the survey shows.

The two panels are \emph{almost} mirror images, and where they differ is
instructive. Offset $+d$ and offset $-d$ are the same permutation exactly when
$2d \equiv 0 \pmod 8$---so offset $\pm 4$ (and the identity) are one operator with
bit-identical panels, whereas offset $\pm 1$ and $\pm 3$ are distinct
\emph{inverse} rotations. Inverses share cycle structure, fixed rules, and the
sustain-or-collapse verdict, but not their transient dynamics: offset $+1$ and
offset $-1$ realise the same sustaining regime through visibly different fields.
\end{workedexample}

\begin{definition}[label={def:braid}]{Temporal braid}
A \emph{temporal braid} alternates two writings between complete coupled spatial
updates: writing $w_A$ is used at even times and $w_B$ at odd times. More
generally, a periodic \emph{schedule} selects one writing at each time step.
This preserves the intermediate spatial states and neighbour blends; it is not
the algebraic composition $P_{\sigma_A}\circ P_{\sigma_B}$ of two simultaneous
maps.
\end{definition}

\begin{definition}[label={def:interrupter}]{Interrupter}
An \emph{interrupter} is an update between elementary propagator writes that is
not itself a propagator write. The coordinate-wise neighbour blend of
Box~\ref{box:metaca-step} is the interrupter in the ordinary spatial runs; a
hold gate that suppresses a scheduled write is another example. An interrupter
changes the coupled dynamics but does not, by definition, guarantee sustained
diversity.
\end{definition}

\begin{proposition}[label={obs:flat}]{Composition parity}
Products alternate between complemented and uncomplemented coordinate permutations.
Each operator of Equation~\eqref{eq:prop} negates every bit
($P_\sig(g)[j] = \lnot\,g[\sig^{-1}(j)]$, a coordinate permutation followed by a
global flip), so under composition the flips cancel in pairs: $P_\sig \circ
P_\tau$ is a \emph{pure} coordinate permutation with no negation left---a
bijection that fixes $2^{c}$ bytes for $c$ cycles but carries none of the forced
$\lambda = 1/2$ structure. An even number of propagators therefore composes to
an uncomplemented coordinate permutation, and an odd number to another member
of the permutation-propagator family.
\end{proposition}

\begin{definition}[label={def:river}]{River}
A \emph{river} is the composed operator obtained by following the collapsing
offset-$+2$ propagator ($\gcd 2$) with a phenotype-reading construction
step---match the local phenotype against four template candidates, first match
wins, and fall back to the all-zero rule when none matches.
\end{definition}

\end{document}
